\chapter{The audit engine: what gets said out loud}
\label{ch:audit}

A classifier produces a label. An auditor produces an accusation. The distance
between those two is the audit engine, and it is where most of the decisions
that a user actually experiences are made.

\section{The decision procedure}

The procedure is ordered, first match wins, and it runs per control:

\begin{enumerate}
  \item If the control is undetectable, emit \code{UNDETECTABLE\_FIELD} with the
    reason and stop.
  \item If a declaration is present and parses to no token, emit
    \code{AUTOCOMPLETE\_OFF} and stop.
  \item If a declaration is present and is not a specification token, emit
    \code{OFF\_SPEC\_TOKEN} naming the token that was meant, and stop.
  \item If a declaration is present and the classifier is confident it names a
    different field, emit \code{WRONG\_AUTOCOMPLETE} and stop.
  \item If no declaration is present and the classifier is confident, emit
    \code{MISSING\_AUTOCOMPLETE} with the token to add.
  \item Otherwise emit the near miss or structural findings that apply, or
    nothing.
\end{enumerate}

Three places in that ordering had to be read rather than transcribed, and each is
a point where two parts of the specification pull against each other.

\textbf{The composite finding has no branch of its own in the published
pseudocode.} Its trigger reads as \emph{an inferred composite}, which taken
literally fires a warning on every correctly declared month and year input,
while the corpus design makes any finding above the lowest severity on a correct
form a defect. Both hold only if the finding is about the \emph{absence} of the
declaration, so it sits where the missing declaration finding sits and carries
the composite's own fix text.

\textbf{The off switch branch has no stop marker in the pseudocode.} Two other
branches carry one and this one does not, but a declaration of \code{off} parses
to no token, so without the stop a field collects both that finding and the
missing declaration finding: two accusations with two contradictory fixes. The
prose above the pseudocode says first match wins, so it stops.

\textbf{The cross origin frame fix text contradicts the extractor's own
guidance.} The generic template ends by suggesting the control be exposed in
light DOM or the attribute declared on the host, which is sound for a closed
shadow root and nonsense for a frame the page does not own. The more specific
rule wins, and the override is data beside the template rather than a branch
inside a renderer.

\section{Thresholds are per engine, and one of them is a mapping}

\begin{table}[htbp]
  \centering
  \caption{One threshold block per engine. An engine with no block is an error
  rather than an inheritance.}
  \label{tab:thresholds}
  \begingroup\small\setlength{\tabcolsep}{4pt}
\begin{tabular}{lrrr>{\raggedright\arraybackslash}p{.36\textwidth}}
\toprule
\textbf{Engine} & \textbf{tau high} & \textbf{tau low} & \textbf{Target} & \textbf{Basis} \\
\midrule
\texttt{rules} & 0.7000 & 0.4000 & null & rule-tier-band-mapping \\  % results: src/autofill_audit/audit/thresholds.json
\texttt{ngram} & 0.8964 & 0.6278 & 0.9800 & dev split of the seed 20260825 corpus, calibrated n-gram engine \\  % results: src/autofill_audit/audit/thresholds.json
\texttt{llm} & 0.0000 & 0.0000 & null & not a boundary: the self-reported scale decides nothing (spec section 12.1) \\  % results: src/autofill_audit/audit/thresholds.json
\texttt{bert} & 0.8813 & 0.4609 & 0.9800 & dev split of the seed 20260825 corpus, calibrated INT8 encoder. Derived under the identical policy of experiments/predictions/p4-threshold-derivation.md, before the test-split run and before the ship condition of spec section 10.7 was evaluated \\  % results: src/autofill_audit/audit/thresholds.json
\bottomrule
\end{tabular}
\endgroup
\par\smallskip\noindent{\footnotesize Source: \rpath{src/autofill_audit/audit/thresholds.json}}

\end{table}

Table~\ref{tab:thresholds} carries a genuine asymmetry and the report should not
smooth it.

The n-gram engine's boundary is derived on the development split against a
precision target that was committed before the model existed. The rule engine's
is a documented mapping from its six tiers onto the same two bands, and every
key in its block that would hold a measurement is null. That is deliberate. A
regular expression table has no probability, and printing a plausible looking
number in a file the runtime reads would be a law 3 violation with a straight
face. The renderers say in words that a rule tier confidence is a tier and not a
frequency.

The language model's block is zero and zero, for the reason given in
Chapter~\ref{ch:classifiers}: a self reported confidence is forbidden from
gating anything, and the engine still needs two numbers.

\begin{table}[htbp]
  \centering
  \caption{The n-gram engine's derived boundary and what it bought on the
  development split. The precision target was registered before the model
  existed and the derivation script was run unmodified.}
  \label{tab:threshold-derivation}
  \begingroup\small\setlength{\tabcolsep}{4pt}
\begin{tabular}{lr}
\toprule
\textbf{Derived quantity} & \textbf{Value} \\
\midrule
Target precision & 0.9800 \\  % results: src/autofill_audit/audit/thresholds.json
tau high & 0.8964 \\  % results: src/autofill_audit/audit/thresholds.json
tau low & 0.6278 \\  % results: src/autofill_audit/audit/thresholds.json
Dev precision at tau high & 0.9954 \\  % results: src/autofill_audit/audit/thresholds.json
Dev recall at tau high & 0.3701 \\  % results: src/autofill_audit/audit/thresholds.json
Dev accusations & 216 \\  % results: src/autofill_audit/audit/thresholds.json
Dev fields needing an accusation & 581 \\  % results: src/autofill_audit/audit/thresholds.json
\bottomrule
\end{tabular}
\endgroup
\par\smallskip\noindent{\footnotesize Source: \rpath{src/autofill_audit/audit/thresholds.json}}

\end{table}

Two things about Table~\ref{tab:threshold-derivation} need saying rather than
leaving in a file.

\textbf{A precision figure is only as strong as its denominator.} A precision of
one over a handful of accusations and a precision of one over a thousand are the
same number and are not the same claim. The derived block therefore records the
denominators beside the rates, so a reader does not have to go and find out which
it is.

\textbf{A high precision target makes a much quieter tool.} At this boundary the
model raises far fewer accusations than there are fields that need one. That is
what a pre-registered high precision target buys on a model whose calibrated
confidences are honest about how often it is right, and it is the correct outcome
of the policy rather than a failure of it. The policy exists because a false
critical costs a developer far more than a missed one, and this is what
\emph{far more} looks like when it is taken literally.

The two engines' boundaries also had to stay separate. The derived n-gram
boundary sits far above the rule engine's confident band, and a single pair of
thresholds across both would have moved every rule engine finding on every page
as a side effect of training a model. The threshold document therefore carries a
block per engine, the loader takes an engine name, and an engine with no block
raises rather than inheriting one. That was designed to bite when the language
model arrived without a block, and it did, which is how the zero and zero
decision came to be made deliberately and in advance.

\section{The finding catalogue, in brief}

\begin{table}[htbp]
  \centering
  \small
  \begin{tabular}{l>{\raggedright\arraybackslash}p{.66\textwidth}}
    \toprule
    \textbf{Severity} & \textbf{Codes} \\
    \midrule
    critical & \ident{MISSING_AUTOCOMPLETE}, \ident{WRONG_AUTOCOMPLETE},
               \ident{OFF_SPEC_TOKEN} \\
    warning  & \ident{AUTOCOMPLETE_OFF}, \ident{UNLABELED_FIELD},
               \ident{PLACEHOLDER_AS_LABEL}, \ident{SPLIT_FIELD},
               \ident{COMPOSITE_FIELD}, \ident{UNDETECTABLE_FIELD} \\
    info     & \ident{GENERIC_IDENTIFIER}, \ident{WRONG_INPUT_TYPE},
               \ident{MISSING_NAME_ATTR}, \ident{EXTRACTION_INCOMPLETE} \\
    note     & \ident{LOW_CONFIDENCE} \\
    \bottomrule
  \end{tabular}
  \caption{The finding catalogue by severity. Each code carries a trigger, a
  named evidence list, a confidence or a documented tier, and a one line fix.}
  \label{tab:findings-catalogue}
\end{table}

Each carries a trigger, a severity, a named evidence list, a confidence or a
documented tier, and a one line fix as text. There is no mode that edits the
user's markup, because rewriting a production template from a probabilistic
classifier's output is exactly the failure the first law exists to prevent.

\section{A worked example}

The following is real output against a checkout fixture authored in this
repository to get things wrong in the ways real checkouts get them wrong.
Nothing is edited except for trimming the middle of the table.

\begin{shell}
$ autofill-audit audit tests/fixtures/checkout_hostile.html
autofill-audit  file:///.../tests/fixtures/checkout_hostile.html

critical (6)

 control          finding
 .................................................................
 #ck-email        MISSING_AUTOCOMPLETE
                  add autocomplete="email" to #ck-email
                  evidence: declaration:absent, label:email-words
                  [rule tier HIGH]
 #ck-postcode     OFF_SPEC_TOKEN
                  autocomplete="zipcode" is not a valid autofill
                  token; use "postal-code"
                  evidence: declaration:off-spec, label:postcode-words
                  [rule tier HIGH]
 #ck-holder       WRONG_AUTOCOMPLETE
                  #ck-holder declares autocomplete="name" but looks
                  like cc-name; change to autocomplete="cc-name"
                  evidence: declaration:token-mismatch,
                  label:cardholder-words [rule tier HIGH]

warning (8)

 control          finding
 .................................................................
 #input7          PLACEHOLDER_AS_LABEL
                  #input7 uses a placeholder as its label; add a
                  real <label>
                  evidence: structure:placeholder-is-the-only-label
 frame[#hosted-pan] UNDETECTABLE_FIELD
                  the control is inside a cross-origin frame, which
                  is how hosted payment fields are built on purpose;
                  this is not necessarily a defect. Audit that
                  document separately
                  evidence: structure:undetectable

$ echo $?
1
\end{shell}

The before and after that a developer actually applies is one attribute:

\begin{markup}
before, nothing tells the browser what this control is for:

  <input type="text" id="ck-email" placeholder="name@example.com">

after, one attribute, and the browser fills it:

  <input type="email" id="ck-email" autocomplete="email"
         placeholder="name@example.com">
\end{markup}

Three properties of that output are the whole design. Every finding names its
evidence, so it can be argued with. The rule engine's confidence is a tier and
prints as a tier, because turning a regular expression table's output into a
percentage would assert a frequency nobody has measured. And a correct field is
met with silence: the correctly declared card number input on that page does not
appear in the report at all.

\section{The exit code contract}

The process exits non zero when a finding at or above the configured failure
threshold is present, which is what lets the tool drop into a pipeline without a
wrapper script. The threshold is configurable, the default is the critical
severity, and the contract is exercised by end to end tests that run the tool in
a subprocess so that the exit codes are the ones a shell sees rather than the
ones a test harness believes.
